\documentclass[UTF8,a4paper,12pt,fontset=none]{ctexart}

\usepackage[a4paper,margin=2.2cm]{geometry}
\usepackage{algorithm}
\usepackage{algpseudocode}
\usepackage{amsmath}
\usepackage{array}
\usepackage{booktabs}
\usepackage{caption}
\usepackage{enumitem}
\usepackage{float}
\usepackage{fontspec}
\usepackage{hyperref}
\usepackage{listings}
\usepackage{siunitx}
\usepackage{tabularx}
\usepackage{tikz}
\usepackage{xcolor}
\usetikzlibrary{arrows.meta,positioning,shapes.geometric}

\hypersetup{colorlinks=true,linkcolor=blue!60!black,urlcolor=blue!60!black}
\defaultCJKfontfeatures{AutoFakeSlant=.2}
\setmainfont[SlantedFont=Times New Roman Italic,BoldItalicFont=Times New Roman Bold Italic]{Times New Roman}
\setCJKmainfont[SlantedFont=Songti SC]{Songti SC}
\setCJKsansfont{Heiti SC}
\setCJKmonofont{Heiti SC}
\setmonofont{Menlo}
\setlength{\parindent}{2em}
\setlength{\parskip}{0.35em}
\captionsetup{font=small}
\sisetup{round-mode=places,round-precision=6}
\renewcommand{\figurename}{图}
\renewcommand{\tablename}{表}
\floatname{algorithm}{算法}
\renewcommand{\lstlistingname}{代码}

\definecolor{CodeBg}{HTML}{F7F7F7}
\definecolor{CodeFrame}{HTML}{DDDDDD}
\definecolor{AccentA}{HTML}{1D4E89}
\definecolor{AccentB}{HTML}{0B6E4F}
\definecolor{AccentC}{HTML}{C73E1D}

\lstdefinestyle{cstyle}{
  language=C,
  basicstyle=\ttfamily\small,
  keywordstyle=\color{AccentA}\bfseries,
  commentstyle=\color{AccentB},
  stringstyle=\color{AccentC},
  numbers=left,
  numberstyle=\tiny\color{gray},
  breaklines=true,
  frame=single,
  rulecolor=\color{CodeFrame},
  backgroundcolor=\color{CodeBg},
  tabsize=2,
  showstringspaces=false
}

\title{\bfseries 中山大学计算机学院本科实验报告}
\author{}
\date{}

\begin{document}

\begin{titlepage}
  \centering
  {\zihao{2}\bfseries 中山大学计算机学院本科实验报告\par}
  \vspace{0.5em}
  {\zihao{4}（2025学年春季学期）\par}
  \vspace{2em}

  \begin{tabularx}{0.9\textwidth}{>{\bfseries}p{4cm}X}
    课程名称： & 并行程序设计与算法 \\
    实验题目： & Lab4 基于 Pthreads 的并行方程求解与蒙特卡洛算法 \\
    学号： & 23336128 \\
    姓名： & 梁力航 \\
    完成日期： & 2026年4月23日 \\
  \end{tabularx}

  \vfill
  \begin{minipage}{0.88\textwidth}
    \small
    本实验使用 Pthreads 完成两个并行程序：其一通过条件变量协调判别式计算线程与两个求根线程，求解一元二次方程；其二将蒙特卡洛采样任务分配给多个线程，使用线程局部随机数种子估计圆周率。实验重点不只是给出正确结果，还包括分析线程创建、条件变量同步、随机数生成与任务粒度对性能的影响。
  \end{minipage}
\end{titlepage}

\begin{abstract}
本实验围绕 Pthreads 编程模型实现两个典型任务：并行求解一元二次方程与蒙特卡洛方法估计圆周率。方程求解任务采用条件变量构造显式任务依赖，使两个求根线程在判别式 \(\Delta=b^2-4ac\) 计算完成后继续执行。蒙特卡洛任务将总采样数 \(n\) 静态划分到多个线程，每个线程使用 \texttt{rand\_r()} 维护独立随机数状态，最后由主线程汇总落入四分之一圆的点数 \(m\)，得到 \(\pi\approx 4m/n\)。实验结果表明，细粒度方程求解任务主要受线程创建和同步开销支配；在题目给定采样规模内，蒙特卡洛任务可以展示一定并行收益，但线程创建与调度开销仍会明显影响加速比。
\end{abstract}

\section{实验目的}

本实验目标如下：
\begin{enumerate}[leftmargin=2.2em]
  \item 使用 Pthreads 条件变量实现线程间同步，完成一元二次方程求根；
  \item 使用 Pthreads 对蒙特卡洛采样任务进行线程级并行划分；
  \item 使用线程安全的随机数生成方式，避免多个线程共享 \texttt{rand()} 全局状态；
  \item 记录不同线程数下的运行时间、估计值与加速比，并分析线程创建和同步开销。
\end{enumerate}

\section{Algorithm Design}

\subsection{并行一元二次方程求解}

一元二次方程 \(ax^2+bx+c=0\) 的根由公式
\[
  x=\frac{-b\pm\sqrt{b^2-4ac}}{2a}
\]
给出。由于两个根都依赖判别式 \(\Delta=b^2-4ac\)，本实验将任务拆为三个线程：
\begin{itemize}
  \item 线程 A 计算判别式 \(\Delta\)，写入共享状态后调用 \texttt{pthread\_cond\_broadcast()}；
  \item 线程 B 等待条件变量，收到通知后计算 \(x_1\)；
  \item 线程 C 等待条件变量，收到通知后计算 \(x_2\)。
\end{itemize}

\begin{algorithm}[H]
  \caption{基于条件变量的一元二次方程并行求解}
  \begin{algorithmic}[1]
    \State 初始化互斥锁、条件变量与共享标志 \(\textit{delta\_ready}\gets false\)
    \State 创建两个求根线程 \(T_1,T_2\)，二者进入条件变量等待
    \State 创建判别式线程 \(T_\Delta\)
    \State \(T_\Delta\) 计算 \(\Delta=b^2-4ac\)
    \State \(T_\Delta\) 加锁，写入 \(\Delta\)，设置 \(\textit{delta\_ready}\gets true\)
    \State \(T_\Delta\) 广播条件变量，唤醒 \(T_1,T_2\)
    \State \(T_1,T_2\) 分别计算 \(\frac{-b+\sqrt{\Delta}}{2a}\) 与 \(\frac{-b-\sqrt{\Delta}}{2a}\)
    \State 主线程 join 三个线程并输出根与运行时间
  \end{algorithmic}
\end{algorithm}

\begin{figure}[H]
  \centering
  \begin{tikzpicture}[
    node distance=1.9cm and 2.3cm,
    task/.style={rectangle,rounded corners,draw=AccentA,thick,minimum width=3.2cm,minimum height=1cm,align=center,fill=blue!4},
    wait/.style={diamond,draw=AccentC,thick,aspect=2.2,align=center,fill=red!4},
    edge/.style={-{Latex[length=2.5mm]},thick}
  ]
    \node[task] (start) {主线程\\创建 Pthreads};
    \node[task,below=of start] (delta) {线程 A\\计算 \(\Delta=b^2-4ac\)};
    \node[wait,below left=of delta] (wait1) {线程 B\\等待 cond};
    \node[wait,below right=of delta] (wait2) {线程 C\\等待 cond};
    \node[task,below=of wait1] (x1) {计算 \(x_1\)};
    \node[task,below=of wait2] (x2) {计算 \(x_2\)};
    \node[task,below right=1.8cm and -0.15cm of x1] (join) {主线程 join\\输出结果};

    \draw[edge] (start) -- (delta);
    \draw[edge] (start) -- (wait1);
    \draw[edge] (start) -- (wait2);
    \draw[edge] (delta) -- node[above left,sloped]{broadcast} (wait1);
    \draw[edge] (delta) -- node[above right,sloped]{broadcast} (wait2);
    \draw[edge] (wait1) -- (x1);
    \draw[edge] (wait2) -- (x2);
    \draw[edge] (x1) -- (join);
    \draw[edge] (x2) -- (join);
  \end{tikzpicture}
  \caption{一元二次方程求解任务依赖图}
  \label{fig:quadratic-dag}
\end{figure}

\subsection{蒙特卡洛方法估计圆周率}

蒙特卡洛方法在单位正方形 \([0,1]\times[0,1]\) 内生成 \(n\) 个随机点。若点满足 \(x^2+y^2\le 1\)，则它位于单位圆的四分之一圆内。设命中点数为 \(m\)，则
\[
  \pi \approx 4\frac{m}{n}.
\]

并行版本将 \(n\) 按线程数 \(p\) 静态切分，每个线程只维护自己的命中计数，最后由主线程汇总。为了避免 \texttt{rand()} 的全局状态竞争，每个线程使用独立 seed 调用 \texttt{rand\_r()}。

\begin{algorithm}[H]
  \caption{Pthreads 蒙特卡洛估计圆周率}
  \begin{algorithmic}[1]
    \State 输入总采样数 \(n\)、线程数 \(p\)、随机种子 \(s\)
    \State 将 \(n\) 划分为 \(p\) 份，余数分配给前若干线程
    \For{每个线程 \(i\)}
      \State 使用独立 seed 初始化线程局部随机数状态
      \State 生成 \((x,y)\in[0,1]^2\)，统计 \(x^2+y^2\le 1\) 的次数 \(m_i\)
    \EndFor
    \State 主线程 join 所有线程，计算 \(m=\sum_i m_i\)
    \State 输出 \(n,m,\pi=4m/n\) 与运行时间
  \end{algorithmic}
\end{algorithm}

\section{Implementation Details}

\subsection{条件变量同步实现}

代码 \ref{lst:cond} 展示了方程求解程序的关键同步逻辑。求根线程使用 \texttt{while} 循环检查共享标志，而不是只用一次 \texttt{if}，这是条件变量的常规写法，可抵御虚假唤醒。判别式线程完成计算后先在互斥锁保护下写入共享状态，再广播条件变量唤醒两个求根线程。

\begin{lstlisting}[style=cstyle,caption={条件变量同步核心代码},label={lst:cond}]
pthread_mutex_lock(&state->mutex);
while (!state->delta_ready) {
    pthread_cond_wait(&state->cond, &state->mutex);
}
const double delta = state->delta;
pthread_mutex_unlock(&state->mutex);

pthread_mutex_lock(&state->mutex);
state->delta = delta;
state->delta_ready = true;
pthread_cond_broadcast(&state->cond);
pthread_mutex_unlock(&state->mutex);
\end{lstlisting}

完整源码位于 \texttt{lab/lab4/src/quadratic\_solver\_pthreads.c}。程序支持浮点数输入，例如：
\begin{lstlisting}[style=cstyle]
./bin/quadratic_solver_pthreads 1 -3 2
\end{lstlisting}
输出保持稳定的 \texttt{key=value} 形式，包括 \texttt{delta}、\texttt{x1}、\texttt{x2} 和 \texttt{time\_sec}。

\subsection{线程局部随机数实现}

代码 \ref{lst:randr} 展示了蒙特卡洛任务的核心循环。每个线程拥有独立 \texttt{seed}，调用 \texttt{rand\_r()} 时只更新自己的局部状态，因此不需要在随机数生成阶段加锁。每个线程也只写入自己的 \texttt{inside} 字段，最终汇总在所有线程 join 之后完成。

\begin{lstlisting}[style=cstyle,caption={蒙特卡洛线程局部采样核心代码},label={lst:randr}]
static void *monte_carlo_worker(void *arg) {
    MonteCarloTask *task = (MonteCarloTask *)arg;
    unsigned int seed = task->seed;
    uint64_t inside = 0;

    for (uint64_t i = 0; i < task->samples; ++i) {
        const double x = unit_random(&seed);
        const double y = unit_random(&seed);
        if (x * x + y * y <= 1.0) {
            ++inside;
        }
    }

    task->inside = inside;
    task->seed = seed;
    return NULL;
}
\end{lstlisting}

完整源码位于 \texttt{lab/lab4/src/monte\_carlo\_pi\_pthreads.c}。程序使用方式如下：
\begin{lstlisting}[style=cstyle]
./bin/monte_carlo_pi_pthreads 65536 8 20250401
\end{lstlisting}

\section{Performance Analysis}

\subsection{实验环境与运行方式}

程序使用如下命令编译。计时使用单调时钟，统计范围覆盖线程创建、线程计算、线程 join 与结果汇总，因此表格反映的是端到端运行时间。

\begin{lstlisting}[style=cstyle]
make -C lab/lab4
./lab/lab4/bin/quadratic_solver_pthreads 1 -3 2
./lab/lab4/bin/monte_carlo_pi_pthreads 65536 8 20250401
\end{lstlisting}

\subsection{实验结果}

表 \ref{tab:mc} 给出了题目输入范围上限 \(n=65,536\) 时蒙特卡洛程序在不同线程数下的单次运行结果。加速比以 1 线程时间为基准计算。

\begin{table}[H]
  \centering
  \caption{蒙特卡洛估计 \(\pi\) 的运行时间与加速比}
  \label{tab:mc}
  \renewcommand{\arraystretch}{1.18}
  \begin{tabular}{rrrrr}
    \toprule
    线程数 & 采样数 \(n\) & 命中数 \(m\) & 时间 \(t\) / s & 加速比 \\
    \midrule
    1 & 65536 & 51445 & 0.000719 & 1.000 \\
    2 & 65536 & 51414 & 0.000384 & 1.872 \\
    4 & 65536 & 51419 & 0.000258 & 2.787 \\
    8 & 65536 & 51524 & 0.000253 & 2.842 \\
    \bottomrule
  \end{tabular}
\end{table}

表 \ref{tab:quadratic} 给出方程求解程序的代表性输出。由于该任务固定由 3 个线程完成，并且单次计算量极小，其运行时间主要由线程创建、条件变量等待和调度开销构成。

\begin{table}[H]
  \centering
  \caption{一元二次方程并行求解代表性结果}
  \label{tab:quadratic}
  \renewcommand{\arraystretch}{1.18}
  \begin{tabularx}{0.92\textwidth}{rrrrlX}
    \toprule
    \(a\) & \(b\) & \(c\) & \(\Delta\) & 根类型 & 输出结果 \\
    \midrule
    1 & -3 & 2 & 1 & 实根 & \(x_1=2,\ x_2=1,\ t\approx 1.50\times 10^{-4}s\) \\
    1 & 2 & 5 & -16 & 复根 & \(x_1=-1+2i,\ x_2=-1-2i,\ t\approx 2.69\times 10^{-4}s\) \\
    \bottomrule
  \end{tabularx}
\end{table}

\subsection{性能讨论}

对一元二次方程求解而言，并行化的教育意义大于性能意义。计算 \(\Delta\) 和两个根都只是常数级浮点运算，而创建 3 个线程、维护互斥锁、进入条件变量等待、广播唤醒和上下文调度都需要额外成本。因此该任务通常不会比串行公式求解更快。它适合作为条件变量同步练习，而不适合作为追求加速比的计算任务。

对蒙特卡洛估计 \(\pi\) 而言，每个采样点之间相互独立，任务具有天然数据并行性。表 \ref{tab:mc} 中 2、4、8 线程相对单线程均有加速，但由于题目给定的最大采样数只有 \(65,536\)，单次计算时间已经进入亚毫秒量级，线程创建、join 和操作系统调度开销占比较高。因此 4 线程到 8 线程的收益明显收敛，实际加速比低于理想线性值。若进一步扩大采样规模，固定线程开销会被更多随机采样计算摊薄；但在本实验题面范围内，更合理的结论是：该任务具备并行结构，实际收益受任务粒度限制。

本实验没有在采样循环中使用全局锁。每个线程只写自己的局部计数，主线程在 join 后汇总，因此同步开销集中在线程生命周期边界。这种设计符合 KISS 和 DRY 原则：并行路径简单、共享状态最小、无需引入复杂的锁竞争控制。

\section{Conclusion}

本实验完成了两个 Pthreads 并行程序和对应报告源码。方程求解任务展示了条件变量的正确使用方式：共享状态必须由互斥锁保护，等待端必须循环检查谓词，生产端写入状态后再广播唤醒。蒙特卡洛任务展示了更适合并行化的数据并行结构：采样任务彼此独立，只需在线程结束后进行一次归约。

实验结论是：并行程序设计不能只看“是否使用多线程”，还要判断任务粒度是否足够大、同步是否必要、共享状态是否可避免。对于极小计算任务，多线程开销可能显著超过计算本身；对于可分解的大规模独立采样任务，Pthreads 可以有效提升吞吐量。

\section*{附录：源码文件}

本实验源码组织如下：
\begin{itemize}
  \item \texttt{lab/lab4/src/quadratic\_solver\_pthreads.c}：条件变量同步的一元二次方程求解；
  \item \texttt{lab/lab4/src/monte\_carlo\_pi\_pthreads.c}：线程局部随机数的蒙特卡洛估计 \(\pi\)；
  \item \texttt{lab/lab4/Makefile}：统一构建入口。
\end{itemize}

提交 PDF 时，可按课程要求命名为 \texttt{并行程序设计\_学号\_姓名.pdf}。

\end{document}
